\chapter[Design of SwarmScan]{Design of SwarmScan, a Cooperative Inventory
    System for Unknown Warehouses}%
\label{chap:system}

\subimport{.}{intro}

\section{Problem Formulation}%
\label{sec:sys:problem}

\subsection{The Inspection Task}%
\label{subsec:sys:task}

An inventory inspection compares what is physically present on the shelves with
the records held by the \gls{abb:wms}. Keeping these two views consistent is what
makes the operation necessary, since miscounts and misplaced items reduce order
accuracy and raise operating costs~\cite{pore2026drone}. The task takes as input
a bounded space containing stored items, each carrying a visual identifier on its
faces, and returns the list of identifiers actually found, together with the
position at which each was read. Two properties determine whether that list is
usable. It must account for as large a share of the stored items as possible,
since an item that is not counted is an item the records will continue to report
incorrectly. And it must contain no identifier that does not exist in the
warehouse, since a single spurious entry places the whole list in doubt.

\subsection{Operating Conditions}%
\label{subsec:sys:conditions}

The task is carried out under conditions that no design can change. They are
stated here as facts, and the requirements of the following section follow from
them.

\textbf{The space}: No description of the building is available at take-off. The
number of racks, their position and the location of the items are unknown, and
the free volume is bounded by fixed obstacles, racks, columns and pallets left on
the floor, as well as by moving ones, operators and handling
vehicles~\cite{pore2026drone, pore2026uavqr}. A typical warehouse is a large
building whose racks can exceed ten metres in
height~\cite{kalinov2020warevision}, and the spacing between them is tight enough
to force conservative speed planning~\cite{pore2026drone}. The contents of the
space also change while the mission runs: items are deposited and removed during
the operating day, and an aisle that was clear a minute earlier can be
obstructed. Since neither the layout nor its evolution is known beforehand, no
route through the space and no division of that space between vehicles can be
established before the flight.

\textbf{The reading}: A visual identifier is decodable only within a bounded range
of distances, within a bounded viewing angle and below a bounded relative
speed~\cite{dong2024motionblurred, alsulami2026deep}, so a surface can be
observed without its label being read. Lighting and occlusion narrow that range
further: decoding accuracy falls from 98.5\% in a well-lit environment to 85.2\%
in low light and to 78.3\% when the code is partly
hidden~\cite{alsulami2026deep}, and the motion of the platform produces blur that
lowers the recognition rate~\cite{dong2024motionblurred}. The total number of
items present is not known before the mission, and nothing in the environment
signals that the inventory is complete.

\textbf{The fleet}: The endurance of a single vehicle is bounded, and falls below
what the coverage of a working warehouse demands: small \glspl{abb:uav} typically
fly for twenty to thirty minutes~\cite{mohsan2023unmanned}, while a complete
manual inventory check consumes more than three days of warehouse
activity~\cite{kalinov2020warevision}. Several vehicles therefore operate in the
same space, which includes aisles wide enough for only one of them at a time.
Using more than one also shortens the mission, since two drones halved the
inspection time of a single one in a warehouse experiment~\cite{pore2026uavqr}. A
vehicle can stop flying at any moment, from a technical failure or a depleted
battery, and without warning~\cite{chandran2024network}. Perception and decision
run on the same limited onboard computer, and the mission depends on both, a
combination identified in the literature as an open
difficulty~\cite{sandino2020autonomous}.

\subsection{Requirements and Constraints}%
\label{subsec:sys:requirements}

From these conditions follow seven results that a solution is expected to
produce. They are stated as outcomes that can be verified on a completed mission
rather than as mechanisms, so that they remain independent of any particular
design.

\begin{enumerate}[itemsep=0.5em]
    \item The inventory is completed without any description of the layout, any
          route, and any division of work being supplied before take-off.
    \item No identifier is reported that does not correspond to an item present
          in the warehouse.
    \item When a vehicle stops flying, the area it was covering is inspected by
          the remaining vehicles before the mission ends.
    \item Two vehicles do not travel to the same objective, so that the capacity
          of the fleet is not spent twice on a single reading.
    \item No vehicle collides with a surface, with an obstacle that appeared
          after take-off, or with another vehicle.
    \item The mission ends at a moment the system determines by itself, without
          an external instruction, and without spending the remaining flight
          budget on a search that no longer produces readings.
    \item The computation performed for each decision, together with the
          perception running alongside it, remains within the onboard budget and
          within the endurance of the platform.
\end{enumerate}

\subsection{Position of Existing Approaches}%
\label{subsec:sys:position}

No existing system produces these seven results together. Systems designed for
warehouse inspection receive the layout before take-off and compute their
coverage route once, on the ground, leaving the vehicle to execute
it~\cite{kalinov2020warevision, pore2026uavqr, cristiani2020inventory,
    alsulami2026deep, kwon2020robust, guinand2021uavugv}. What matters is not how
that route is computed but when: the decision is made when the system knows the
least it will ever know, and a target appearing later cannot be taken into
account, because the component that produced the plan is no longer
running~\cite{krayani2023goal}. The same design carries two further assumptions.
The environment is treated as immobile, although operators and handling vehicles
move through a working warehouse, and only about 3\% of inventory studies
consider a dynamic one~\cite{pore2026drone}. And cooperation between flying
vehicles is almost absent: where several vehicles are used, the ground robot
carries the drone or places beacons rather than reading any code
itself~\cite{kalinov2020warevision, guinand2021uavugv}.

The methods that could decide during the flight are not designed for this task.
Those that coordinate several vehicles fix the division of labour before
take-off, by a central schedule~\cite{guinand2021uavugv}, a fixed
split~\cite{pore2026uavqr}, or a leader broadcasting
waypoints~\cite{chandran2024network}, and allocation able to adapt in real time
is still described as a future direction rather than a current
practice~\cite{kumar2025comprehensive}. Semantic reasoning appears in two
separate places, and never where it would be needed. On a ground robot working
alone, it does drive exploration, but the mission solves Object Goal Navigation
and ends as soon as one instance of a designated target is
reached~\cite{yokoyama2024vlfm, yu2023l3mvn}, and the representation built for
that target cannot be reused for another~\cite{yokoyama2024vlfm}. On a drone, it
only interprets an instruction given by a person: the vehicle carries out what it
has been asked to do and never chooses where to explore~\cite{chen2024typefly}.
Semantic reasoning, cooperation and flight therefore never occur together; three
of the four combinations exist, and the missing one is exactly what an inventory
requires.

The two sides therefore never meet: the systems that perform the right task
cannot react, and the systems that can react were not built for it and stop at
the first success. Four of the seven results are ruled out by that alone. An
inventory cannot be completed without a layout when the layout is what the
planner consumes. The area of a vehicle that stops cannot be reassigned when the
assignment was made once. An obstacle appearing after take-off cannot be avoided
by a route fixed before it existed. And a system that ends at the first detection
cannot decide when an inventory is complete. One caveat applies throughout: these
limitations are established almost entirely on simulated missions, since about
45\% of inventory studies are simulation-based against roughly 7\% that report
real experiments~\cite{pore2026drone}.

\section{System Overview}%
\label{sec:sys:overview}

\section{The Shared Map}%
\label{sec:sys:map}

\section{Two-Track Perception}%
\label{sec:sys:perception}

\section{Generation and Scoring of Targets}%
\label{sec:sys:targets}

\section{Coordination Mechanism}%
\label{sec:sys:coordination}

\section{Semantic Guide}%
\label{sec:sys:semantic}

\section{Flight Controller}%
\label{sec:sys:control}

\section{A Complete Mission}%
\label{sec:sys:mission}

\section{Approaches Explored and Set Aside}%
\label{sec:sys:discarded}
